%% kapa: exam=1 round=29 year=2018 period=2 form=A subject=appraisal_law subject_ko=감정평가관계법규 seq=1 number=1
%% source: Q-Net 감정평가사 기출문제 (cst00309) artlSeq=5202808 file=extracted/1008/5202808/2018년 제29회 감정평가사 1차시험 기출문제/2018년 제29회 감정평가사 1차시험 기출문제/2교시 A형.pdf
%% conversion: kordoc 4.24.51 -> tools/gen1.py (auto); review=proofread-sonnet; flags=none
\begin{question}{1}
국토의 계획 및 이용에 관한 법령상 도시ㆍ군기본계획에 관한 설명으로 옳은 것은?
\begin{choices}
  \item 시장 또는 군수는 도시ㆍ군기본계획의 수립을 위한 공청회 개최와 관련한 사항을 일간신문에 공청회 개최예정일 7일 전까지 2회 이상 공고하여야 한다.
  \item 도시ㆍ군기본계획에는 기후변화 대응 및 에너지절약에 관한 사항에 대한 정책 방향이 포함되어야 한다.
  \item 시장 또는 군수는 3년마다 관할 구역의 도시ㆍ군기본계획에 대하여 그 타당성 여부를 전반적으로 재검토하여 정비하여야 한다.
  \item 시장 또는 군수가 도시ㆍ군기본계획을 변경하려면 지방의회의 승인을 받아야 한다.
  \item 시장 또는 군수는 대통령령이 정하는 바에 따라 도시ㆍ군기본계획의 수립기준을 정한다.
\end{choices}
\end{question}
\answer{2}
